\chapter{Enlarged Prostate (Benign Prostatic Hyperplasia)}

\section*{Definition}
An enlarged prostate (benign prostatic hyperplasia, BPH) is a noncancerous increase in prostate size that commonly affects older men. BPH can compress the urethra and cause bothersome lower urinary tract symptoms.

\section*{What Happens in Your Body}
BPH involves hyperplasia of stromal and glandular elements in the transition zone of the prostate. As the prostate enlarges, it compresses the prostatic urethra, increasing resistance to urine flow. Bladder detrusor muscle hypertrophies to overcome the obstruction, leading to instability, reduced compliance, and eventual decompensation with incomplete emptying.

\section*{Causes}
\begin{itemize}
\item Aging (BPH affects over 50 percent of men by age 60 and 90 percent by age 85)
\item Androgenic hormones (dihydrotestosterone drives prostatic growth)
\item Genetic and familial factors
\item Ethnic variation (higher in African-American men)
\item Obesity
\item Lack of physical activity
\item Erectile dysfunction medications (may worsen symptoms)
\end{itemize}

\section*{When to See a Doctor}
\begin{itemize}
\item Frequent urination (urinary frequency, especially at night)
\item Urgency to urinate
\item Weak urine stream or difficulty starting urination
\item Dribbling at the end of urination
\item Feeling of incomplete bladder emptying
\item Straining to urinate
\item Sudden inability to urinate (acute urinary retention)
\end{itemize}

\section*{Related Symptoms}
\begin{itemize}
\item Frequent urination
\item Nocturia (nighttime urination)
\item Urinary urgency
\item Overactive bladder
\item Urinary retention (inability to urinate)
\item Urinary tract infection (from incomplete emptying)
\item Hematuria (blood in urine)
\item Prostate cancer
\end{itemize}

\section*{Self-Care / Home Management}
\begin{itemize}
\item Limit fluid intake before bedtime
\item Reduce caffeine and alcohol consumption
\item Avoid decongestants and antihistamines
\item Empty the bladder completely when urinating
\item Double-void (wait a moment then try again)
\item Maintain a healthy weight
\item Stay physically active
\item Practice timed voiding (urinate on schedule)
\end{itemize}

\section*{How Your Doctor Will Evaluate You}
\begin{itemize}
\item Digital rectal examination
\item International Prostate Symptom Score questionnaire
\item Urinalysis and urine culture
\item Serum PSA (prostate-specific antigen)
\item Urinary flow rate measurement (uroflowmetry)
\item Post-void residual volume measurement
\item Transrectal ultrasound of the prostate
\end{itemize}

\section*{Treatment Options}
\begin{itemize}
\item Watchful waiting for mild symptoms
\item Alpha-blockers (tamsulosin, terazosin) for symptom relief
\item 5-alpha reductase inhibitors (finasteride, dutasteride) to shrink the prostate
\item Combination therapy for moderate-to-severe symptoms
\item TURP (transurethral resection of the prostate)
\item Minimally invasive procedures (Rezum, UroLift)
\item Surgery for severe cases or complications
\end{itemize}
\section*{References}
\begin{thebibliography}{99}
\bibitem{ref1} McVary KT, Roehrborn CG, Avins AL, et al. "Update on AUA Guideline on the Management of Benign Prostatic Hyperplasia." J Urol. 2011;185(5):1793--1803.
\bibitem{ref2} Berry SJ, Coffey DS, Walsh PC, et al. "The Development of Human Benign Prostatic Hyperplasia with Age." J Urol. 1984;132(3):474--479.
\bibitem{ref3} Parsons JK, Dahm P, Kohler TS, et al. "AUA Guideline: Surgical Management of Lower Urinary Tract Symptoms Attributed to Benign Prostatic Hyperplasia." J Urol. 2020;204(4):799--804.
\bibitem{ref4} Thorpe A, Neal D. "Benign Prostatic Hyperplasia." Lancet. 2003;361(9366):1359--1367.
\bibitem{ref5} Lepor H. "Medical Treatment of Benign Prostatic Hyperplasia." Rev Urol. 2011;13(1):20--28.
\bibitem{ref6} Madersbacher S, Alivizatos G, Nordling J, et al. "EAU 2004 Guidelines on Assessment, Therapy and Follow-Up of Men with Lower Urinary Tract Symptoms Suggestive of Benign Prostatic Obstruction." Eur Urol. 2004;46(5):547--554.
\end{thebibliography}


\clearpage
